%%%%%%%%%%%%%%%%%
% This is an example CV created using altacv.cls (v1.1.5, 1 December 2018) written by
% LianTze Lim (umairkhan@softtechconsultants.com), based on the
% Cv created by softtech consultants at http://www.softtechconsultants.com
%
%% It may be distributed and/or modified under the
%% conditions of the LaTeX Project Public License, either version 1.3
%% of this license or (at your option) any later version.
%% The latest version of this license is in
%%    http://www.latex-project.org/lppl.txt
%% and version 1.3 or later is part of all distributions of LaTeX
%% version 2003/12/01 or later.
%%%%%%%%%%%%%%%%

%% If you are using \orcid or academicons
%% icons, make sure you have the academicons
%% option here, and compile with XeLaTeX
%% or LuaLaTeX.
% \documentclass[10pt,a4paper,academicons]{altacv}

%% Use the "normalphoto" option if you want a normal photo instead of cropped to a circle
% \documentclass[10pt,a4paper,normalphoto]{altacv}

\documentclass[10pt,a4paper,ragged2e]{altacv}

%% AltaCV uses the fontawesome and academicon fonts
%% and packages.
%% See texdoc.net/pkg/fontawecome and http://texdoc.net/pkg/academicons for full list of symbols. You MUST compile with XeLaTeX or LuaLaTeX if you want to use academicons.

% Change the page layout if you need to
\geometry{left=2cm,right=10cm,marginparwidth=6.8cm,marginparsep=1.2cm,top=1.25cm,bottom=1.25cm}

% Change the font if you want to, depending on whether
% you're using pdflatex or xelatex/lualatex
\ifxetexorluatex
% If using xelatex or lualatex:
\setmainfont{Arial}
\else
% If using pdflatex:
\usepackage[utf8]{inputenc}
\usepackage[T1]{fontenc}
\usepackage[default]{lato}
\fi

% Change the colours if you want to
\definecolor{VividPurple}{HTML}{000000}
\definecolor{SlateGrey}{HTML}{2E2E2E}
\definecolor{LightGrey}{HTML}{2E2E2E}
\colorlet{heading}{VividPurple}
\colorlet{accent}{VividPurple}
\colorlet{emphasis}{SlateGrey}
\colorlet{body}{LightGrey}

% Change the bullets for itemize and rating marker
% for \cvskill if you want to
\renewcommand{\itemmarker}{{\small\textbullet}}
  \renewcommand{\ratingmarker}{\faCircle}

  %% sample.bib contains your publications
  \addbibresource{sample.bib}

  \begin{document}
  \name{Tomas Xavier Santos}
  \tagline{QA Engineer | Software PM}
  \personalinfo{%
    % Not all of these are required!
    % You can add your own with \printinfo{symbol}{detail}
    \email{tom.xaviersantos@gmail.com}
    \phone{+55 81 998975432}
    %  \mailaddress{Address, Street, 00000 County}
    \location{Recife, Brazil}
    %  \homepage{marissamayr.tumblr.com/}
    %  \twitter{@marissamayer}
    \linkedin{www.linkedin.com/in/tomasxs/}
    % I'm just making this up though.
    %  \orcid{orcid.org/0000-0000-0000-0000} % Obviously making this up too. If you want to use this field (and also other academicons symbols), add "academicons" option to \documentclass{altacv}
    }

    %% Make the header extend all the way to the right, if you want.
    \begin{fullwidth}
      \makecvheader
    \end{fullwidth}

    %% Depending on your tastes, you may want to make fonts of itemize environments slightly smaller
    \AtBeginEnvironment{itemize}{\small}

    %% Provide the file name containing the sidebar contents as an optional parameter to \cvsection.
    %% You can always just use \marginpar{...} if you do
    %% not need to align the top of the contents to any
    %% \cvsection title in the "main" bar.
    \cvsection[page1sidebar]{Description}
    QA Software Engineer \& Software Project Manager with 4+ years of specialized experience in Android \& Web manual testing, test automation, and quality assurance. Proficient in automating processes and optimizing workflows, including AI-driven products. Reduced localization testing cycles by 90\% through custom automation frameworks, and successfully delivered 30+ mobile projects for Motorola, currently managing 10+ active projects simultaneously for Motorola. Proven track record of delivering high-quality software solutions in fast-paced environments.

    %% Divides

    \cvsection[]{Experience}

    \cvevent{Project Manager}{Roundtable Studio - Contractor for Motorola Mobility }{Dec 2024 -- Present}{Recife, Brazil}
    \begin{itemize}
      \item Orchestrated end-to-end management of 30+ mobile projects, driving project delivery, timelines, and cross-functional collaboration between Marketing, Development, and Localization teams.
      \item Established and delivered a comprehensive training program for a team of 8 new QA engineers, incorporating onboarding documentation and conducting a week long hands-on program for localization related Android testing best practices.
      \item Proactively managed project risks related to software readiness and resource availability, implementing contingency plans that maintained 100\% on-time delivery across all projects.
    \end{itemize}

    \divider

    \cvevent{QA Engineer}{Projeto/Cin Motorola} {Jun 2022 -- Nov 2024}{Recife, Brazil}
    \begin{itemize}
      \item Spearheaded Android QA for mobile products, owning the complete testing lifecycle from planning to production release and ensuring zero critical production bugs.
      \item Developed automated testing framework using Python and API automation to ensure translation quality for indigenous endangered languages, reducing testing time from 10+ hours to less than 1 hour (90\% reduction) while eliminating human errors.
      \item Drove and took ownership of quality initiatives for AI-based products, designing specialized test strategies related to localization for ML-powered features to ensure robust testing coverage.
    \end{itemize}

    \divider

    \cvevent{QA Tester - Intern}{Projeto/Cin Motorola} {Aug 2021 -- May 2022}{Recife, Brazil}
    \begin{itemize} 
      \item Validated UI elements on Android, Desktop and Web applications across 50+ languages, verifying text truncation, right-to-left language support, character encoding, and culturally appropriate imagery to ensure compliance with Motorola's localization standards.
    \end{itemize}

    %\divider

    % \divider
    %\cvskill{German}{3}


    % \divider

    % \cvevent{Product Engineer}{Google}{23 June 1999 -- 2001}{Palo Alto, CA}

    % \begin{itemize}
    % \item Joined the company as employe \#20 and female employee \#1
    % \item Developed targeted advertisement in order to use user's search queries and show them related ads
    % \end{itemize}

    %\cvsection{A Day of My Life}

    % Adapted from @Jake's answer from http://tex.stackexchange.com/a/82729/226
    % \wheelchart{outer radius}{inner radius}{
    % comma-separated list of value/text width/color/detail}
    % Some ad-hoc tweaking to adjust the labels so that they don't overlap
    % \wheelchart{1.5cm}{0.5cm}{%
    %   10/10em/accent!30/Sleeping \& dreaming about work,
    %   25/9em/accent!60/Public resolving issues with Yahoo!\ investors,
    %   5/13em/accent!10/\footnotesize\\[1ex]New York \& San Francisco Ballet Jawbone board member,
    %   20/15em/accent!40/Spending time with family,
    %   5/8em/accent!20/\footnotesize Business development for Yahoo!\ after the Verizon acquisition,
    %   30/9em/accent/Showing Yahoo!\ employees that their work has meaning,
    %   5/8em/accent!20/Baking cupcakes
    % }

    \clearpage

    %\cvsection[page1sidebar]{Publications}

    \nocite{*}

    % \printbibliography[heading=pubtype,title={\printinfo{\faBook}{Books}},type=book]

    % \divider

    % \printbibliography[heading=pubtype,title={\printinfo{\faFileTextO}{Journal Articles}}, type=article]

    % \divider

    % \printbibliography[heading=pubtype,title={\printinfo{\faGroup}{Conference Proceedings}},type=inproceedings]

    % %% If the NEXT page doesn't start with a \cvsection but you'd
    % %% still like to add a sidebar, then use this command on THIS
    % %% page to add it. The optional argument lets you pull up the
    % %% sidebar a bit so that it looks aligned with the top of the
    % %% main column.
    % % \addnextpagesidebar[-1ex]{page3sidebar}


  \end{document}
